% 用法: xelatex SL_gap_n1_inf_limit_proof.tex (两次, -output-directory=build)
% 主题: 定理 A (相邻间距 INF R->inf 极限) 的 T1/T2/T3 严格证明;
%       证明见 §2 (数值仅作证据, 见 §3 计算机辅助认证, §4 数值证据 (非证明)).
\documentclass[UTF8]{ctexart}
\usepackage{amsmath, amssymb, amsthm}
\usepackage[margin=2.5cm]{geometry}
\usepackage[hyphens]{url}
\usepackage[hidelinks]{hyperref}
\usepackage{booktabs}
\usepackage{enumitem}
\emergencystretch 3em

\newtheorem{theorem}{定理}[section]
\newtheorem{proposition}[theorem]{命题}
\newtheorem{lemma}[theorem]{引理}
\newtheorem{corollary}[theorem]{推论}
\newtheorem{remark}[theorem]{注}
\newtheorem{conjecture}[theorem]{猜想}

\title{SL 相邻特征间距的 $R\to\infty$ 极限: 第八轮修订证明}
\author{研究组 (BVE research, run R-20260806T200000Z-inflimit-5B2C7D, 会话 30)}
\date{2026-09-22 修订; 原稿 2026-08-07}

\begin{document}
\maketitle

\begin{abstract}
考虑 Dirichlet 问题 $-y''=\lambda\rho y$ ($\rho\in[1,R]$ a.e.) 中对称三块
$\rho=R$ on $[0,u]\cup[1-u,1]$, $\rho=1$ on $(u,1-u)$ 的相邻特征间距
$D_R(u)=\lambda_2-\lambda_1$ 的下确界:
\begin{equation}
	\lim_{R\to\infty} R\cdot m_R \;=\; \bar D(u^*)\;=\; 24.943866138432\ldots
	\;<\; 3\pi^2,
	\qquad m_R:=\inf_{u\in(0,1/2)} D_R(u),
\end{equation}
其中 $u^*=0.32992250812006654959\ldots$ 是 $S(u)=0$ 在 $(0,1/2)$ 的唯一解
(等价地 $\bar D'(u^*)=0$), 即 $\bar D$ 的全局极小点 $u^*$.
证明分三步: (T1) 极限存在, (T2) $\bar D$ 的严格单调结构, (T3) 区间包含.
(T2) 通过符号链 (K~$\to J\to G\to S$) 初等地建立;
(T1) 的核心是引理 A$''$: 当 $w:=u\sqrt R\ge 2$ 时 $R\cdot D_R(u)\ge\bar D(u)$,
其初等修订给出更强正余量 $G-\bar D>\ell/(Ru^3)$.
薄层区域 ($w\le2$) 用连续相位速度的统一下界覆盖, 得 $G>15\pi^2/4>3\pi^2$.
T1 因而不依赖旧二维网格或 T3 的高精度数值; T3 由新的精确有理证书支撑.
正文结构: §2 严格证明, §3 计算机辅助认证, §4 数值证据 (非证明).
\end{abstract}

\tableofcontents

\section{问题与记号}

\subsection{设定}
考虑 Dirichlet 问题
\begin{equation}
	-y''(x)=\lambda\,\rho(x)\,y(x),\qquad y(0)=y(1)=0,
\end{equation}
对称三块权: 对 $u\in(0,1/2)$,
\begin{equation}
	\rho_{R,u}(x)=R\ \text{on }(0,u)\cup(1-u,1),\qquad
	\rho_{R,u}(x)=1\ \text{on }(u,1-u).
\end{equation}
$0<\lambda_1(R,u)<\lambda_2(R,u)$ 为前两个特征值,
\begin{equation}
	D_R(u):=\lambda_2(R,u)-\lambda_1(R,u),\qquad
	m_R:=\inf_{u\in(0,1/2)}D_R(u),\qquad
	G(R,u):=R\cdot D_R(u).
\end{equation}
记 $\varepsilon:=R^{-1/2}$, $\ell:=1/2-u$, $w:=u\sqrt R=u/\varepsilon$.
用缩放特征值 $\mu_k(R,u):=R\lambda_k(R,u)$, 则 $G(R,u)=\mu_2-\mu_1$.

\subsection{极限系统}
对 $u\in(0,1/2)$ 令 $\bar\theta_2=\bar\theta_2(u)=a(u)\in(\pi/2,\pi)$ 为
\begin{equation}
	\tan a=a\Bigl(1-\frac1{2u}\Bigr)
\end{equation}
的唯一根 (等价地 $\tan a=-a\ell/u$), 并令
\begin{equation}
	\bar\mu_1(u):=\frac{\pi^2}{4u^2},\qquad
	\bar\mu_2(u):=\frac{a(u)^2}{u^2},\qquad
	\bar D(u):=\bar\mu_2(u)-\bar\mu_1(u).
\end{equation}
以及
\begin{equation}
	S(u):=\bar\mu_1(u)\frac2u-\bar\mu_2(u)\frac{\sin^2 a(u)}{I_2(u)},
	\qquad I_2(u):=\int_0^u\sin^2\bigl(\sqrt{\bar\mu_2(u)}\,x\bigr)\,dx
	=\frac u2-\frac{u\sin 2a(u)}{4a(u)}.
\end{equation}

\subsection{主定理}
\begin{theorem}[定理 A (相邻间距的 INF 极限)]
设 $u^*\in(0,1/2)$ 为 $S(u)=0$ 的唯一根 (等价地 $\bar D'(u^*)=0$). 则
\begin{equation}
	\lim_{R\to\infty} R\,m_R=\bar D(u^*),\qquad
	\bar D(u^*)=24.943866138432\ldots<3\pi^2=29.608813203268\ldots,
\end{equation}
并且任何近极小化子序列 $u_R$ (即 $D_R(u_R)\le m_R+R^{-2}$) 满足 $u_R\to u^*$.
\end{theorem}

\begin{remark}[证明结构]
定理 A 的证明分解为三个定理:
\begin{itemize}
	\item \textbf{T1} (收敛): $\lim_R R\,m_R=\bar D(u^*)$ 及近极小化子收敛到 $u^*$ (见 §\ref{sec:T1});
	\item \textbf{T2} (单调结构): $S$ 在 $(0,1/2)$ 有唯一零点 $u^*$, 且 $\bar D$ 在 $u^*$ 处全局严格极小 (见 §\ref{sec:T2});
	\item \textbf{T3} (区间包含): $\bar D(u^*)$ 的高精度区间包含且 $\bar D(u^*)<3\pi^2$ (见 §\ref{sec:T3}).
\end{itemize}
其余为初等引理; 这里的 $\inf$ 是固定 $R$ 下的 $\inf$, 全局对称性问题属 O3a/C1, 另行处理.
\end{remark}

\section{严格证明}

证明先确定正确的模式分支. 大 $w$ 区域用初等差量界 $d_1>4\alpha$, $d_2<3\alpha$;
薄层区域用连续展开相位的速度上界, 统一得到 $G>15\pi^2/4$.
T1 只需这些连续估计和 T2 的解析单调结构; T3 的精确数值包含单独证明.

\subsection{全域匹配方程与有条件的相位分支}\label{sec:global-phase}
令 $\varepsilon=R^{-1/2}$, $\ell=1/2-u$, $\theta_k=u\sqrt{R\lambda_k}$,
$z_k=\ell\sqrt{\lambda_k}$, $c=\varepsilon\ell/u$.
对称性、特征值单性及振荡定理说明第一模态为偶函数, 第二模态为奇函数.
第二模态若为偶函数, 则其零点成对, 而中心同时满足函数值与导数为零又与唯一性矛盾;
这与第二模态恰有一个内部零点不符. 两个模态限制到 $(0,1/2)$ 后均可取为正.
重层解为 $A\sin(\sqrt{R\lambda}\,x)$; 轻层的偶、奇解分别为
$B\cos(\sqrt\lambda(1/2-x))$ 与 $B\sin(\sqrt\lambda(1/2-x))$.
在 $x=u$ 匹配函数和导数, 消去振幅, 得到全域无极点方程
\begin{align}
 E(\theta,z)&:=\cos\theta\cos z-\varepsilon\sin\theta\sin z=0
 &&\text{(偶)},\label{eq:polefree-even}\\
 O(\theta,z)&:=\varepsilon\sin\theta\cos z+\cos\theta\sin z=0
 &&\text{(奇)}.\label{eq:polefree-odd}
\end{align}
这里尚未除以任何三角因子. 半区间基态正性给出
$0<\theta_k<\pi$, $0<z_1<\pi/2$ 及 $0<z_2<\pi$.
偶模在接口处的轻层导数为正, 所以 $\cos\theta_1>0$,
进而 $0<\theta_1<\pi/2$ 对全部 $R>1$, $0<u<1/2$ 成立.
奇模的 $\theta_2>\pi/2$ 则不能无条件使用.

以下相位括号与引理 A$''$ 明确假设 $R\ge1500$, $w=u/\varepsilon\ge2$.
此时 $0<c=1/(2w)-\varepsilon<1/4$, 故 $z_2=c\theta_2<\pi/4$.
由 \eqref{eq:polefree-odd} 可知 $\cos\theta_2<0$, 即
$\pi/2<\theta_2<\pi$. 相关分母均非零, 因而可在此范围写成
\begin{equation}
 \cot\theta_1=\varepsilon\tan z_1,\qquad
 \tan\theta_2=-\varepsilon^{-1}\tan z_2.\label{eq:secular}
\end{equation}
令 $\delta_1=\pi/2-\theta_1$, $\delta_2=\theta_2-\pi/2$. 此范围内有
\begin{equation}
 \tan\delta_1=\varepsilon\tan z_1,\quad
 z_1=(\pi/2-\delta_1)c;\qquad
 \tan\delta_2=\varepsilon\cot z_2,\quad
 z_2=(\pi/2+\delta_2)c.\label{eq:delta}
\end{equation}
从而
\begin{equation}
 \mu_1=\frac{(\pi/2-\delta_1)^2}{u^2},\quad
 \mu_2=\frac{(\pi/2+\delta_2)^2}{u^2},\quad
 G(R,u)=\frac{(\pi/2+\delta_2)^2-(\pi/2-\delta_1)^2}{u^2}.\label{eq:mu}
\end{equation}
例如 $R=1600$, $u=1/1600$ 时, $\rho\ge1$ 与 min-max 给出
$\lambda_2\le4\pi^2$, 因此 $\theta_2\le\pi/20<\pi/2$.
这个实例反驳全域奇相位分支的旧说法, 不反驳上述受限分支.

\subsection{相位括号 (引理 2.1)}
\begin{lemma}\label{lem:phases}
设 $R\ge 1500$, $u\ge 2\varepsilon$ (即 $w\ge 2$). 则
\begin{enumerate}
	\item[(a)] $0\le\delta_1\le\delta_1^+$, 其中
		$\delta_1^+:=\arctan\bigl(\varepsilon\tan\bigl(\frac\pi2\frac{\varepsilon\ell}{u}\bigr)\bigr)
		\le\varepsilon\tan\frac\pi8<0.011<\frac\pi4$;
	\item[(b)] $0\le\delta_2\le\delta_2^+$, 其中 $\delta_2^+:=\arctan\frac{2u}{\pi\ell}$;
	\item[(c)] $z_2\le\pi/8$;
	\item[(d)] $\psi_2:=\bar\theta_2-\theta_2\ge 0$, 即 $\theta_2\le\bar\theta_2=a(u)$.
\end{enumerate}
\end{lemma}
\begin{proof}
由 $u\ge2\varepsilon$ 得 $\ell/u\le1/(4\varepsilon)$, 从而
$\varepsilon\ell/u\le1/4$ 且 $\alpha:=(\ell/u)\varepsilon^2\le\varepsilon/4$.

(a) $z_1=(\pi/2-\delta_1)\varepsilon\ell/u\le(\pi/2)\varepsilon\ell/u\le\pi/8$,
$\tan z_1$ 递增, 故 $\delta_1=\arctan(\varepsilon\tan z_1)\le\delta_1^+$;
再用 $\arctan x\le x$ 与 $\tan(\pi/8)<0.415$, 得
$\delta_1^+\le\varepsilon\tan(\pi/8)\le\varepsilon_0\tan(\pi/8)<0.011<\pi/4$.

(b) 对 $0<z<\pi$ 有 $\cot z\le1/z$; 因 $z_2=\theta_2\varepsilon\ell/u<\pi\cdot(1/4)<\pi$, 故
$\tan\delta_2=\varepsilon\cot z_2\le\varepsilon/z_2=u/(\theta_2\ell)\le u/((\pi/2)\ell)$,
故 $\delta_2\le\arctan(2u/(\pi\ell))=\delta_2^+$.

(c) $z_2=(\pi/2+\delta_2)\varepsilon\ell/u\le(\pi/2+\delta_2^+)\cdot x$, $x:=\varepsilon\ell/u$.
考察 $h(x):=x\bigl(\pi/2+\arctan(2\varepsilon/(\pi x))\bigr)$ 在 $(0,1/4-\varepsilon]$ 上单调:
\begin{equation}
	h'(x)=\frac\pi2+\arctan\frac{2\varepsilon}{\pi x}-\frac{2\varepsilon\pi x}{\pi^2x^2+4\varepsilon^2}
	\ge\frac\pi2-\frac12>0,
\end{equation}
其中 $\sup_{x>0}2\varepsilon\pi x/(\pi^2x^2+4\varepsilon^2)=1/2$ (在 $x=2\varepsilon/\pi$ 处取得).
故 $h(x)\le h(1/4-\varepsilon)$, 而
\begin{align}
	h\Bigl(\frac14-\varepsilon\Bigr)
	&=\Bigl(\frac14-\varepsilon\Bigr)
	\Bigl(\frac\pi2+\arctan\frac{8\varepsilon}{\pi(1-4\varepsilon)}\Bigr)
	\le\Bigl(\frac14-\varepsilon\Bigr)
	\Bigl(\frac\pi2+\frac{2\pi\varepsilon}{1-4\varepsilon}\Bigr)\notag\\
	&=\frac\pi8-\frac{\pi\varepsilon}{2}
	+\frac{2\pi\varepsilon(1/4-\varepsilon)}{1-4\varepsilon}
	=\frac\pi8-\frac{\pi\varepsilon}{2}+\frac{\pi\varepsilon}{2}
	=\frac\pi8,
\end{align}
其中用到 $\arctan t\le t$ 与 $8/\pi\le2\pi$ (因 $4\le\pi^2$).

(d) 令 $g(\theta):=\theta-\pi/2-\arctan(\varepsilon\cot(\theta\varepsilon\ell/u))$.
$g$ 在 $(\pi/2,\pi)$ 上严格递增 ($g'>0$), 且 $\theta_2$ 是 $g$ 的零点 (见 \eqref{eq:delta}).
又 $\bar z_2:=\bar\theta_2\varepsilon\ell/u\le\pi/4<\pi$. 由 $\cot z\le1/z$:
\begin{equation}
	\varepsilon\cot\bar z_2\le\frac{\varepsilon}{\bar z_2}
	=\frac{u}{\bar\theta_2\ell}=\tan\bar d_2,\qquad \bar d_2:=\bar\theta_2-\frac\pi2,
\end{equation}
故 $g(\bar\theta_2)=\bar d_2-\arctan(\varepsilon\cot\bar z_2)\ge0$, 且 $g(\bar\theta_2)\ge0=g(\theta_2)$.
由 $g$ 递增得 $\theta_2\le\bar\theta_2$. \qed
\end{proof}

\subsection{引理 A\texorpdfstring{$''$}{} 的初等修订}\label{sec:lemAdp}
\begin{theorem}[引理 A$''$, 带正余量]\label{thm:lemAdp}
设 $R\ge1500$, $0<u<1/2$, $w=u\sqrt R\ge2$. 则
\begin{equation}
 G(R,u)>\bar D(u)+\frac{\ell}{Ru^3}>\bar D(u).
\end{equation}
\end{theorem}
\begin{proof}
令 $c=\varepsilon\ell/u$, $\alpha=\varepsilon c=\ell/(Ru)$,
$t=a(u)$, $\theta=\theta_2$, $v=u/\ell=-t\cot t>0$.
由前节的模式识别, $0<\theta_1<\pi/2<\theta_2<\pi$, 所有下述三角除法合法.
还有 $0<c\le1/4-\varepsilon$, $\varepsilon<1/38$, $\alpha<1/152$.
奇方程与 $\cot z<1/z$ 给出
\[
 \delta_2\le\varepsilon\cot z_2<\frac{\varepsilon}{c\theta_2}
 \le\frac{2\varepsilon}{\pi c},\qquad
 z_2<\frac\pi8-\varepsilon(\pi/2-2/\pi)<\frac\pi8.
\]
函数 $s\mapsto-s\cot s$ 在 $(\pi/2,\pi)$ 上严格递增, 而
$-\theta\cot\theta=\varepsilon\theta\cot(c\theta)<\varepsilon/c=v=-t\cot t$,
所以 $\theta<t$. 定义
\[
 d_1=\pi^2/4-\theta_1^2>0,\qquad d_2=t^2-\theta^2>0,\qquad
 G-\bar D=(d_1-d_2)/u^2.
\]

\emph{第一差量下界.} $z_1<\pi/8$ 且 $\tan(\pi/8)=\sqrt2-1<1/2$,
故 $0<\delta_1<1/76$ 及 $\theta_1>7/5$.
用 $\tan z\ge z$, $\arctan x\ge x-x^3/3$, $\pi^2<10$ 得
\[
 \delta_1\ge\arctan(\alpha\theta_1)
 \ge\alpha\theta_1(1-\pi^2\alpha^2/12)
 >\frac{99}{100}\alpha\theta_1.
\]
因此
\begin{equation}
 d_1=(\pi/2+\theta_1)\delta_1
 >\frac{29}{10}\frac75\frac{99}{100}\alpha
 =\frac{20097}{5000}\alpha>4\alpha.\label{eq:r8-def1}
\end{equation}

\emph{余切余项与一个连续三角界.} 令 $r(z)=1/z-\cot z$.
由 $\sin z-z\cos z=\int_0^z s\sin s\,ds$,
$\sin s\le s$ 与 $\sin z\ge z-z^3/6$, 对 $0<z\le\pi/8$ 有
\begin{equation}
 0<r(z)\le\frac{z}{3(1-z^2/6)}
 <\frac{64}{187}z<\frac7{20}z.\label{eq:r8-cot-elementary}
\end{equation}
这里仅用 $3<\pi<22/7$, 因而 $\pi^2<10$.
另对 $\pi/2<t<\pi$,
\begin{equation}
 B(t):=\frac{2t^4}{t^2+v^2+v}
 =\frac{2t^3\sin^2t}{t-\sin t\cos t}<2(t\sin t)^2\le8.\label{eq:r8-B8}
\end{equation}
最后一步可直接证明: $t\le2$ 时 $t\sin t\le2$.
$t\ge2$ 时 $q(t)=t\sin t$ 满足 $q''=2\cos t-t\sin t<0$.
交错 Taylor 余项给出
\[
 \sin2\le\frac{2578}{2835}<\frac{91}{100},\qquad
 \cos2\le-\frac{131}{315}<-\frac25.
\]
故凹函数的切线估计与 $t-2<8/7$ 给出
\[
 q(t)\le q(2)+q'(2)(t-2)
 <\frac{91}{50}+\frac{11}{100}\frac87
 =\frac{681}{350}<2.
\]

\emph{第二差量上界.} 写 $\psi=t-\theta\in(0,\pi/2)$,
$A=\tan(t-\pi/2)=v/t$, $B=\tan(\theta-\pi/2)=\varepsilon\cot(c\theta)$.
恒等式
\[
 A-B=-\frac{v\psi}{t\theta}+\varepsilon r(c\theta),\qquad
 \tan\psi=\frac{A-B}{1+AB}
\]
结合 $\tan\psi\ge\psi$ 给出
\[
 \psi\le\frac{\varepsilon r(c\theta)}{1+AB+v/(t\theta)}.
\]
取 $C=7/20$, 由 \eqref{eq:r8-cot-elementary} 有
\[
 \varepsilon r(c\theta)<C\alpha\theta,\qquad
 AB>\frac{v^2}{t\theta}-C\varepsilon^2\frac\theta t.
\]
令 $D_0=1+v(v+1)/(t\theta)\ge1$, $d=C\varepsilon^2\theta/t<1/1000$.
固定 $t,v$ 时, $\theta(t+\theta)/D_0$ 的正分子随 $\theta$ 增加,
正分母随 $\theta$ 减少, 因而整个比值增加. 由 $\theta<t$ 和 \eqref{eq:r8-B8},
\begin{equation}
 d_2=(t+\theta)\psi
 < C\alpha\frac{\theta(t+\theta)}{D_0-d}
 <\frac7{20}\frac{1000}{999}\,8\alpha
 =\frac{2800}{999}\alpha<3\alpha.\label{eq:r8-def2}
\end{equation}
所有分母均有上述正下界. 由 \eqref{eq:r8-def1} 和 \eqref{eq:r8-def2},
$G-\bar D>\alpha/u^2=\ell/(Ru^3)$, 得证.
\end{proof}

本修订主链仅需上述初等粗界; 特别地 $d_2/d_1<3/4<0.8256$.
旧主链所用更细的 $C_z<0.337$ 和 $B(t)\le9$ 另经第八轮精确证书修复,
作为可复用的标量工具保留, 不再承担 T1 的连续区域覆盖.
\subsection{整个薄层区域的解析下界}\label{sec:sliver}
\begin{lemma}[连续相位速度给出的谱隙下界]\label{lem:r8-phase-speed}
对任意 $R\ge1$, $0<u<1/2$, 令 $\varepsilon=R^{-1/2}$,
$\ell=1/2-u$, $w=u/\varepsilon$. 则
\begin{equation}
 G(R,u)\ge\frac{\pi^2}{2\varepsilon(w+\ell)(w+\varepsilon\ell)}.
 \label{eq:r8-global-gap-bound}
\end{equation}
本结论不要求第二个重层相位大于 $\pi/2$.
\end{lemma}
\begin{proof}
对 $k=\sqrt\lambda>0$, 取 $y(0)=0$, $y'(0)=1$ 的射击解.
接口向量为 $(y'(u),ky(u))=(\cos(wk),\varepsilon\sin(wk))$.
经过轻层后, 它旋转了角 $\ell k$, 两分量即 \eqref{eq:polefree-even}--\eqref{eq:polefree-odd}.
令 $A_\varepsilon$ 是向量 $\cos\theta+i\varepsilon\sin\theta$ 的连续展开辐角,
$A_\varepsilon(0)=0$. 向量从不为零, 且
\[
 A_\varepsilon'(\theta)=\frac{\varepsilon}{\cos^2\theta+\varepsilon^2\sin^2\theta}>0,
 \quad A_\varepsilon(\pi/2)=\pi/2,\quad A_\varepsilon(\pi)=\pi.
\]
半区间末端相位为
\[
 \Phi(k)=\ell k+A_\varepsilon(wk),\qquad
 0<\Phi'(k)\le\ell+w/\varepsilon.
\]
它从零严格增加到无穷. 第一 Neumann 根和第一 Dirichlet 根分别满足
$\Phi(k_1)=\pi/2$, $\Phi(k_2)=\pi$.
空间相位在两层中均严格递增, 故这两个半区间解内部为正.
前者偶延拓没有内部零点, 后者奇延拓恰有中点这一个零点;
由振荡定理, $k_1^2,k_2^2$ 确为全区间的前两个特征值.

对这两个实际根积分相位导数, 得
\[
 k_2-k_1\ge\frac{\pi}{2(\ell+w/\varepsilon)}.
\]
又 $\Phi(\pi/(2w))>\pi/2$, 故 $wk_1<\pi/2$.
在第一分支, 因 $0<\varepsilon\le1$,
$A_\varepsilon(\theta)=\arctan(\varepsilon\tan\theta)\le\theta$,
所以 $k_1\ge\pi/[2(\ell+w)]$.
相乘即得
\[
 G=\varepsilon^{-2}(k_2-k_1)(k_2+k_1)
 \ge\frac{\pi^2}{2\varepsilon(w+\ell)(w+\varepsilon\ell)}.
\]
以上分母均正, 没有使用任何全域奇相位分支假设.
\end{proof}

\begin{lemma}[薄层区域, 替代旧网格证书]\label{lem:sliver}
对每个 $R\ge1500$ 及 $0<u\le2/\sqrt R$, 有
\begin{equation}
 G(R,u)>\frac{15\pi^2}{4}>25,\qquad G(R,u)>3\pi^2.
 \label{eq:r8-sliver-margin}
\end{equation}
\end{lemma}
\begin{proof}
此时 $w\le2$, $\ell<1/2$; 由 $38^2<1500$ 得 $\varepsilon<1/38$.
故 \eqref{eq:r8-global-gap-bound} 给出
\[
 G>\frac{\pi^2}{2(1/38)(5/2)(2+1/76)}
 =\frac{2888}{765}\pi^2>\frac{15}{4}\pi^2.
\]
最后一项有理比较等价于 $11552>11475$.
再用 $\pi>3$, 得 $15\pi^2/4>135/4>25$; 另有 $15/4>3$.
这覆盖整个连续区域, 包括任意小的正 $w$, $w=1/2$, 两条旧曲边界及无穷大的 $R$ 尾部.
\end{proof}

\subsection{T2: \texorpdfstring{$\bar D$}{D} 的单调结构与唯一临界点}\label{sec:T2}
\begin{theorem}[T2]\label{thm:T2}
$S(u)=0$ 在 $(0,1/2)$ 有唯一根 $u^*$, 且 $\bar D$ 在 $u^*$ 处全局严格极小.
\end{theorem}
\begin{proof}
\emph{第 1 步 (参数化).} 令
\begin{equation}
	a\in(\pi/2,\pi)\;\longmapsto\;u(a):=\frac{a}{2(a-\tan a)}\in(0,1/2)
\end{equation}
良定义: 对 $a\in(\pi/2,\pi)$, $\tan a<0$, 故 $a-\tan a>0$, $u(a)>0$,
且 $u(a)\to0^+$ (当 $a\to\pi/2^+$), $u(a)\to1/2^-$ (当 $a\to\pi^-$);
另外
\begin{equation}
	u'(a)=\frac{a-\sin(2a)/2}{2\cos^2a\,(a-\tan a)^2}>0,
\end{equation}
其中 $a-\sin(2a)/2\ge a-1/2>0$ (因 $a\ge\pi/2$).

\emph{第 2 步 (符号链).} 定义
\begin{align}
	\widetilde K(a)&:=-a^2+3\sin^2a+\frac32a\sin2a, &
	J(a)&:=4a^3\cot a+6a^2-\pi^2, &\hspace{-0.6em}
	G(a)&:=8a^3\sin^2a-\pi^2(2a-\sin2a).
\end{align}
符号验证 (初等求导, 也见脚本 07):
\begin{equation}
	J'(a)=\frac{4a\,\widetilde K(a)}{\sin^2a},\qquad
	G'(a)=4\sin^2a\,J(a),
	\label{eq:chain}
\end{equation}
以及
\begin{equation}
	\bar D'(u(a))=S(u(a)),
	\qquad
	S(u(a))=-\frac{4(a-\tan a)^3}{a^3(2a-\sin2a)}\,G(a).
	\label{eq:SG}
\end{equation}

\emph{第 3 步 ($\widetilde K$ 的符号).} 对 $a\in(\pi/2,\pi)$ 有
$\widetilde K(a)=\sin^2a\,h(a)$, $h(a):=3+3a\cot a-a^2/\sin^2a$, 且
\begin{equation}
	h'(a)\sin^3a=3\cos a\,\sin^2a-5a\sin a+2a^2\cos a<0,
\end{equation}
其中 $\cos a<0$, $\sin a>0$, $a>0$ 使三项均为负. 故 $h$ 严格递减; 又
$h(\pi/2)=3-\pi^2/4>0$, $h(a)\to-\infty$ (当 $a\to\pi^-$), 故 $h$ (即 $\widetilde K$)
在 $(\pi/2,\pi)$ 有唯一零点 $a_1$, 且 $\widetilde K>0$ on $(\pi/2,a_1)$,
$\widetilde K<0$ on $(a_1,\pi)$.

\emph{第 4 步 ($J$ 的符号).} 由 \eqref{eq:chain}, $J'=4a\widetilde K/\sin^2a$:
$J$ 在 $(\pi/2,a_1)$ 上递增, 在 $(a_1,\pi)$ 上递减. $J(\pi/2)=\pi^2/2>0$,
$J(a)\to-\infty$ (当 $a\to\pi^-$, 因 $\cot a\to-\infty$). 故 $J$ 有唯一零点
$a^*\in(a_1,\pi)$, $J>0$ on $(\pi/2,a^*)$, $J<0$ on $(a^*,\pi)$.

\emph{第 5 步 ($G$ 的符号).} 由 $G'=4\sin^2a\,J$: $G$ 在 $(\pi/2,a^*)$ 上递增,
在 $(a^*,\pi)$ 上递减. $G(\pi/2)=0$ 且 $G'(\pi/2)=2\pi^2>0$, 故 $G>0$ on $(\pi/2,a^*]$;
$G(\pi)=-2\pi^3<0$. 故 $G$ 在 $(a^*,\pi)$ 内有唯一零点 $a_G$, 即 $G$ 在 $(\pi/2,\pi)$ 上
有唯一零点 $a_G$, 且 $G>0$ on $(\pi/2,a_G)$, $G<0$ on $(a_G,\pi)$.

\emph{第 6 步 ($S$ 的符号).} 由 \eqref{eq:SG}, 因子 $-(4(a-\tan a)^3)/(a^3(2a-\sin2a))<0$
(各因子均为正), 故 $\operatorname{sgn}S(u(a))=-\operatorname{sgn}G(a)$. 于是
$S<0$ on $(0,u^*)$, $S=0$ at $u^*:=u(a_G)$, $S>0$ on $(u^*,1/2)$,
即 $S$ 由负变正. 因 $\bar D'=S$, $\bar D$ 在 $(0,u^*)$ 上递减, 在 $(u^*,1/2)$ 上递增,
故 $u^*$ 为全局严格极小. 为核对端点, 当 $u\downarrow0$ 时写 $a=\pi/2+d$.
根方程等价于 $a\tan d=u/\ell$, 因而 $d/u\to4/\pi$,
$u\bar D(u)=(a+\pi/2)d/u\to4$. 故 $\bar D(u)\to+\infty$.
当 $u\uparrow1/2$ 时 $a\to\pi$, 直接代入得 $\bar D(u)\to3\pi^2$.
严格递增的右支还给出 $\bar D(u^*)<3\pi^2$. \qed
\end{proof}

\begin{remark}[解析唯一性与数值定位]
T2 的唯一性来自上述符号链. T3 只负责以精确包络定位已经确定唯一的根及其函数值.
\end{remark}

\subsection{T3: 区间包含}\label{sec:T3}
\begin{theorem}[T3]\label{thm:T3}
设 $u^*$ 如 T2. 则
\begin{equation}
	u^*\in[0.32992250812006654958,\;0.32992250812006654960],
\end{equation}
\begin{equation}
	\bar D(u^*)\in[24.9438661384324768968,\;24.9438661384324769084],
	\qquad
	3\pi^2-\bar D(u^*)\ge4.664947,
\end{equation}
故 $25-\bar D(u^*)>0.0561$.
\end{theorem}
\begin{proof}[证明 (精确有理包络)]
使用 Machin 恒等式
$\pi=16\arctan(1/5)-4\arctan(1/239)$ 及交错部分和得到 $\pi$ 的有理上下界.
对有理 $x$ 用 $\sin x$, $\cos x$ 的 Taylor 多项式及
$|R_N(x)|\le |x|^{N+1}/(N+1)!$ 作包络; 除法前检查分母区间不含零.
在两个精确有理端点
\begin{align*}
 a_L&=2.27651323902643307501655540074525185589672,\\
 a_H&=2.27651323902643307501655540074525185589673
\end{align*}
分别证 $G(a_L)>0$, $G(a_H)<0$ (此处 $G$ 为 T2 的一变量符号链函数).
由 T2 的唯一性得 $a_G\in(a_L,a_H)$. 使用严格递增的
$u(a)=a/[2(a-\tan a)]$ 传播区间时, 取下端像的下界和上端像的上界,
再对 $(a_G^2-\pi^2/4)/u(a_G)^2$ 使用精确区间四则运算.
所得更窄包络严格位于本定理的两个显示区间内; 比较 $25$ 和 $3\pi^2$
也在有理算术中完成. 十进制显示按向下/向上舍入, 不参与真值判定.
第八轮证书及输入身份见 §\ref{sec:cac}.
旧脚本 05 的两端均取根下界规则不能保证像包含, 不再作为本定理证书.
\end{proof}

\subsection{固定内部参数的解析展开}\label{sec:fixed-u}
\begin{proposition}\label{prop:fixed-u}
固定 $0<u<1/2$, 令 $\ell=1/2-u$, $b=\ell/u$, $a=a(u)$.
当 $R\to\infty$ 时,
\begin{equation}
 G(R,u)=\bar D(u)+\frac{C(u)}R+O_u(R^{-2}),\qquad
 C(u)=\frac{\pi^2b}{2u^2}-\frac{2a^4b^3}{3u^2(1+b+b^2a^2)}.\label{eq:rate}
\end{equation}
余项在任意紧区间 $J\Subset(0,1/2)$ 上可取为 $O_J(R^{-2})$.
在 T2 的极小点,
\begin{equation}
 C(u^*)=\frac{\pi^2(1/2-u^*)}{3(u^*)^3}>0.\label{eq:stationary-rate}
\end{equation}
\end{proposition}
\begin{proof}
令 $t=1/R$. 把无极点偶方程以及除以正数 $\sqrt t$ 的奇方程写成
\begin{align*}
 \mathcal E(\theta,t)&=\cos\theta\cos(b\theta\sqrt t)
 -\sqrt t\sin\theta\sin(b\theta\sqrt t),\\
 \mathcal O(\theta,t)&=\sin\theta\cos(b\theta\sqrt t)
 +\cos\theta\frac{\sin(b\theta\sqrt t)}{\sqrt t}.
\end{align*}
其中 $\cos(b\theta\sqrt t)$, $\sqrt t\sin(b\theta\sqrt t)$,
$\sin(b\theta\sqrt t)/\sqrt t$ 均用幂级数在 $t=0$ 解析延拓.
令 $p=\pi/2$, $H=1+b+b^2a^2>0$. 在 $(p,0)$,
$\mathcal E_\theta=-1$, $\mathcal E_t=-bp$;
在 $(a,0)$, 由 $\sin a+ba\cos a=0$ 得
\[
 \mathcal O_\theta=\cos a\,H\ne0,\qquad
 \mathcal O_t=\frac{b^3a^3\cos a}{3}.
\]
解析隐函数定理于是给出
\[
 \theta_1(t)=p-pbt+O_u(t^2),\qquad
 \theta_2(t)=a-\frac{a^3b^3}{3H}t+O_u(t^2).
\]
对充分小的正 $t$, 相位分别处于 $(0,\pi/2)$ 和 $(\pi/2,\pi)$,
轻层相位足够小, 拼接解在半区间无内部零点. 因此这些根确为所需的偶、奇基态.
平方相减再除以 $u^2$ 得 \eqref{eq:rate}.
在任意 $J\Subset(0,1/2)$ 上, 限制根的导数连续且非零;
有限子覆盖及根的唯一性给出共同邻域和二阶导数界, 从而余项在 $J$ 上一致.

又 $\tan a=-ba$ 给出
\[
 \frac{\sin^2a}{I_2(u)}=\frac{2a^2b^2}{u(1+b+b^2a^2)},\qquad
 I_2(u)=\frac u2(1+b\cos^2a).
\]
代入 $S(u)=2\bar\mu_1/u-\bar\mu_2\sin^2a/I_2$ 可得
\[
 C(u)=\frac{4\bar\mu_1(u)\ell}{3u}+\frac\ell3S(u).
\]
最后用 $S(u^*)=0$ 和 $\bar\mu_1=\pi^2/(4u^2)$ 得 \eqref{eq:stationary-rate}.
\end{proof}
这里不声称 $u\to0$ 或 $u\to1/2$ 时的一致展开, 也不凭固定参数展开交换下确界与极限.
特别地 $\sqrt R(G(R,u^*)-\bar D(u^*))\to0$;
$R(G(R,u^*)-\bar D(u^*))\to C(u^*)$, 其数值约为 $15.5807908501$.

\subsection{T1: 极限、近极小化子与最优值误差}\label{sec:T1}
\begin{theorem}[T1, 解析修订]\label{thm:T1}
设 $u^*$ 为 T2 的唯一极小点, $M=\bar D(u^*)$. 则
\[
 \lim_{R\to\infty}Rm_R=M,\qquad 0\le Rm_R-M=O(R^{-1}).
\]
更一般地, 若 $R_j\to\infty$, $0<u_j<1/2$, $\eta_j\ge0$,
$D_{R_j}(u_j)\le m_{R_j}+\eta_j$ 且 $R_j\eta_j\to0$,
则 $u_j\to u^*$. 原条件 $\eta_j=R_j^{-2}$ 是其中一个特例.
\end{theorem}
\begin{proof}
T2 的严格单调结构及端点极限给出 $M<3\pi^2$.
这一步不使用 T3 的高精度区间. 对每个 $R\ge1500$ 和 $0<u<1/2$,
若 $w\le2$, 引理 \ref{lem:sliver} 给出 $G>15\pi^2/4>3\pi^2>M$;
若 $w\ge2$, 引理 A$''$ 给出 $G\ge\bar D(u)\ge M$.
故 $Rm_R=\inf_u G(R,u)\ge M$.

另一方面, 在单个固定点 $u^*$, 命题 \ref{prop:fixed-u} 或其证明中的相位括号
给出 $G(R,u^*)=M+O(R^{-1})$. 由 $Rm_R\le G(R,u^*)$, 得到极限与非负误差的 $O(R^{-1})$ 上界.
这里用到了前一段的全域下界, 并非仅凭点态展开交换下确界与极限.
本结论没有给出最优值误差的精确首项系数, 也没有给出极小化参数的收敛速度.

对所述近极小化子,
\[
 R_jm_{R_j}\le G(R_j,u_j)\le R_jm_{R_j}+R_j\eta_j\longrightarrow M.
\]
薄层下界的严格余量迫使 $w_j=u_j\sqrt{R_j}>2$ 最终成立, 因而
$M\le\bar D(u_j)\le G(R_j,u_j)\to M$.
任取在 $[0,1/2]$ 中收敛的子列: 极限为零将导致 $\bar D(u_j)\to+\infty$,
与有限极限 $M$ 矛盾; 极限为 $1/2$ 将导致 $\bar D(u_j)\to3\pi^2>M$, 也矛盾.
所以聚点都在内部. 连续性及 T2 的唯一性给出该聚点必须为 $u^*$.
因此整个序列收敛到 $u^*$.
\end{proof}

\begin{corollary}[定理 A]
由 \ref{thm:T1}, \ref{thm:T2}, \ref{thm:T3} 即得定理 A.
\end{corollary}

\section{精确证书、连续覆盖与历史程序}\label{sec:cac}

当前标量证书位于第八轮工件目录中的 \texttt{certificate/certificate.py};
以 Python 标准库 Fraction 作证明算术, 以 Decimal 的定向舍入作十进制显示.
输入有理数、Machin 的交错余项、Taylor 的全局 Lagrange 余项及每次除法的非零分母均明确检查.
实现与独立执行证据见
\href{../reports/proof-audit-round8-20260922/REPORT.md}{第八轮修订报告}.

T3 使用根的异号端点和 T2 的唯一性, 经正确方向的单调像传播得到包含.
新的引理 A$''$ 主链已经解析证明 $B(t)\le8$ 及 $d_2/d_1<3/4$.
旧链所需的 $C_z<0.337$, $B(t)\le9$ 与比例 $<0.8256$ 也由新证书独立支持,
作为可复用标量工具保留; 这些细常数不再是 T1 的必要依赖.
薄层区域的连续覆盖以 §\ref{sec:sliver} 的新证明为准, 不再调用旧矩形网格.

历史 run \texttt{R-20260806T200000Z-inflimit-5B2C7D} 的源文件和输出保持原字节:
\begin{itemize}
 \item 脚本 05 把根的下界误作右端像上界, 且有普通浮点中间运算. 新证书替代其 T3 作用.
 \item 脚本 16 的普通超越函数值加一次 \texttt{nextafter} 不保证包含;
 B/D 曲边区域的矩形内缩、$w\downarrow0$ 和大 $R$ 尾部亦未完整覆盖.
 加密原网格不能消除这种方向错误.
 \item 脚本 19 同样存在余切漏包, 并把未包络的 $v$ 值当区间端点,
 以小于真实 $\pi$ 的浮点数结束网格. 标量结论现由精确有理证书及解析覆盖支撑.
 \item 脚本 17、18 及旧点抽样保留为历史记录, 本轮没有给它们重新授予区间认证或全域证明地位.
\end{itemize}
局部 Lean 的编译、实际声明和公理依赖另有独立证据; 它不等同于把本篇所有谱理论、
隐函数和连续区域估计全部形式化. 各证据的精确边界与失败记录均列于第八轮报告.

\subsection{余切细常数的正确来源}
对于 $0<z<\pi$, NIST DLMF4.22.3 给出
\[
 \frac{1/z-\cot z}{z}=2\sum_{j=1}^{\infty}\frac1{j^2\pi^2-z^2}.
\]
级数在 $(-\pi,\pi)$ 的每个紧子区间上一致收敛, 每项在正半轴严格递增.
故 $(1/z-\cot z)/z$ 严格递增, 并对 $0<z\le\pi/8$ 有
\[
 1/z-\cot z\le C_z z,\qquad
 C_z=\frac{8/\pi-(1+\sqrt2)}{\pi/8}<\frac{337}{1000}.
\]
Laurent 展开中的正确幂次为
$\cot z=1/z-\sum_{k\ge1}c_kz^{2k-1}$,
$c_k=2^{2k}|B_{2k}|/(2k)!>0$, 从 $z/3$ 开始.
不能把紧子区间上的一致收敛说成在整个 $(0,\pi)$ 上一致收敛.
证书用精确 $\pi$, $\sqrt2$ 有理包络核对细常数; 十进制打印值不作为零宽区间输入.

\section{数值证据与收敛阶的范围}\label{sec:numeric}

历史浮点扫描和拟合保存在原 run 中, 本轮不以它们证明连续区域的下界.
第八轮精确证书重新支持 T3 的显示区间; 附带高精度数值只用于检验解析式的行为.

命题 \ref{prop:fixed-u} 已解析证明固定内部 $u$ 的 $R^{-1}$ 展开.
在 $u^*$ 处, 非零系数约为 $15.5807908501$, 因而
$R(G(R,u^*)-\bar D(u^*))$ 趋于这一常数,
而 $\sqrt R(G(R,u^*)-\bar D(u^*))\to0$.
原文由数值误差声称非零 $R^{-1/2}$ 首项, 已撤回.
本节不推断极小化参数或任意近极小化子序列的收敛速度.

\section{涉及到的数学知识}
\begin{itemize}
	\item \textbf{Sturm--Liouville 谱理论}: Dirichlet 问题, 对称/反对称模态,
		转移矩阵, 相位 (secular) 方程.
	\item \textbf{Feynman--Hellmann 定理}: 特征值的参数单调性 (经典);
		脚本 17 中验证 $d\mu_k/du<0$, $d\mu_k/dR>0$.
	\item \textbf{相位括号方法}: 用单调函数夹逼特征相位 (见 \ref{lem:phases}(d)).
	\item \textbf{三角级数余项}: $\tan$, $\cot$ 的级数余项符号与单调性,
		给出显式常数 (见 §2.3).
	\item \textbf{符号链}: $\widetilde K\to J\to G\to S$ 的逐层符号分析
		(见 \ref{thm:T2}).
	\item \textbf{精确包络}: 有理区间四则运算, Machin 恒等式, Taylor 余项及方向正确的像传播.
	\item \textbf{极限系统端点}: 由 $a(u)$ 的极限方程, $u\to1/2$ 时 $\bar D\to3\pi^2$;
		$u\to0$ 时 $\bar D\to+\infty$.
\end{itemize}

\section{文献}
\sloppy
\begin{enumerate}
\raggedright
	\item Keller: The minimum ratio of two eigenvalues, SIAM J.\ Appl.\ Math.\
		\textbf{29}(2) (1975) 279--283,
		\href{https://doi.org/10.1137/0129024}{DOI 10.1137/0129024}.
	\item Mahar, Willner: An extremal eigenvalue problem, Comm.\ Pure Appl.\ Math.\
		\textbf{29}(5) (1976) 517--529,
		\href{https://doi.org/10.1002/cpa.3160290505}{DOI 10.1002/cpa.3160290505}.
	\item Willner, Mahar: An extremal eigenvalue problem II, SIAM J.\ Math.\ Anal.\
		\textbf{13}(4) (1982) 633--643,
		\href{https://doi.org/10.1137/0513040}{DOI 10.1137/0513040}.
	\item NIST Digital Library of Mathematical Functions, \href{https://dlmf.nist.gov/4.22.E3}{公式 4.22.3},
版本 1.2.8 (2026-09-15), 本轮实际读取该公式及极点排除条件.
\item 本项目内部 run (见 \texttt{runs/rigorous-open-math-research/}):
		R-20260806T200000Z-inflimit-5B2C7D (INF 极限, 本文),
		R-20260805T000000Z-gapn1-a1b2c3 (O1/O2/O3b 前期工作),
		R-20260806T151000Z-o1reaudit-5A1C3D (O1 独立复审),
		R-20260806T140000Z-o3ac1-42F931 (O3a C1).
\end{enumerate}

\end{document}